\section{The Children}\label{sec:children}

The joint job loss and the household's response determine what reaches the children: a family income that stays lower
for years, a higher risk that the parents separate, and, in the month of the shock, a household with no earner. This
section reports what the PSID can say about the children's adult outcomes. The answer is a set of associations that
bound, rather than identify, the effect of a father's job loss; the model uses them as bounds, together with
administrative estimates.

\paragraph{Associations.} I compare children whose father figure was displaced at a given age with children whose
father figure was not, controlling for parental income before the displacement, the child's sex and birth year
(Table~\ref{tab:displacement_children}). A displacement before age six is associated with a household income rank 3.9
percentiles lower and 0.28 fewer years of schooling; a displacement at 12--17 with adult earnings about 14\% lower, a
rank 2.5 percentiles lower and 0.29 fewer years of schooling. For White children most of these associations are
statistically significant; for Black children they are significant mainly for displacements at 12--17, and the
estimates are imprecise. Adding family income after the displacement reduces the associations by about a third,
while adding the change in the mother's hours changes nothing; both are outcomes of the job loss, so this split is
descriptive.\footnote{\texttt{IGE/output/father\_displacement.md}; code \texttt{IGE/20.PaperDisplacement.py}.}

\begin{table}[htbp]\centering\small
\caption{The father's displacement and children's adult outcomes}\label{tab:displacement_children}
\resizebox{\textwidth}{!}{%
\begin{tabular}{lcccc}
\toprule
 & \begin{tabular}[b]{@{}c@{}}(1) Log\\lifetime earnings\end{tabular} & \begin{tabular}[b]{@{}c@{}}(2) Log lifetime\\family income\end{tabular} & \begin{tabular}[b]{@{}c@{}}(3) Family\\income rank\end{tabular} & \begin{tabular}[b]{@{}c@{}}(4) Years\\of schooling\end{tabular} \\
\midrule
\multicolumn{5}{l}{\textit{A. All children}} \\
Father displaced at 0--5 & $-0.058$ & $-0.086^{***}$ & $-3.91^{***}$ & $-0.279^{***}$ \\
 & (0.040) & (0.030) & (1.15) & (0.072) \\
Father displaced at 6--11 & $-0.075$ & $-0.057^{*}$ & $-3.63^{***}$ & $-0.312^{***}$ \\
 & (0.046) & (0.031) & (1.24) & (0.079) \\
Father displaced at 12--17 & $-0.154^{***}$ & $-0.065^{**}$ & $-2.51^{**}$ & $-0.294^{***}$ \\
 & (0.053) & (0.028) & (1.04) & (0.076) \\
\multicolumn{5}{l}{\textit{B. Black children}} \\
Father displaced at 0--5 & $-0.001$ & $-0.057$ & $-2.54$ & $-0.117$ \\
 & (0.071) & (0.056) & (1.87) & (0.104) \\
Father displaced at 6--11 & $-0.048$ & $-0.003$ & $-2.07$ & $-0.041$ \\
 & (0.087) & (0.059) & (2.17) & (0.124) \\
Father displaced at 12--17 & $-0.189^{*}$ & $-0.090$ & $-3.35^{*}$ & $-0.317^{***}$ \\
 & (0.101) & (0.061) & (1.98) & (0.122) \\
\multicolumn{5}{l}{\textit{C. White children}} \\
Father displaced at 0--5 & $-0.084^{*}$ & $-0.101^{***}$ & $-4.48^{***}$ & $-0.364^{***}$ \\
 & (0.047) & (0.035) & (1.44) & (0.095) \\
Father displaced at 6--11 & $-0.102^{*}$ & $-0.089^{**}$ & $-4.68^{***}$ & $-0.511^{***}$ \\
 & (0.054) & (0.037) & (1.52) & (0.102) \\
Father displaced at 12--17 & $-0.137^{**}$ & $-0.057^{*}$ & $-2.47^{**}$ & $-0.324^{***}$ \\
 & (0.062) & (0.030) & (1.22) & (0.098) \\
\bottomrule
\end{tabular}}
\begin{minipage}{0.95\textwidth}\footnotesize Each coefficient is from a separate regression of the outcome in the column on an indicator that the child's father figure (the male head of the child's family) had his first displacement when the child was in the age band of the row, against children whose father figure was not displaced; children first exposed in another band are dropped. Controls: log parental income before the displacement, child sex, Black (pooled panel) and birth-year fixed effects. Adult outcomes are means over ages 25--55 (ranks within birth cohort). Associations, not causal effects. Standard errors clustered by 1968 family. *** $p<0.01$, ** $p<0.05$, * $p<0.1$.\end{minipage}
\end{table}


\paragraph{How much is the job loss.} Displaced families differ before the job loss: their family income is 24 log
points lower, the parents have 0.6 fewer years of schooling, and the father's earnings were already growing more slowly
(all $p<0.01$). Controlling for the father's earnings path, occupation and industry barely changes the associations, but
four checks narrow what can be attributed to the job loss. With plant closures alone the associations vanish for Black
children (adult earnings $+0.085$, s.e.\ 0.078) and survive for White children's schooling ($-0.40$ years, s.e.\ 0.11)
and earnings rank ($-3.3$ percentiles, s.e.\ 1.6). A displacement after the child has left childhood (at 18--25) is not
associated with schooling ($-0.08$, s.e.\ 0.09), while one at 12--17 is ($-0.27$, s.e.\ 0.08; difference $p=0.060$).
The behavior problems of younger children are elevated before the displacement, so that association reflects the
families rather than the job loss. Quits also predict less schooling (Table~\ref{tab:displacement_definitions}), which
signals selection. The main definition reproduces the association of about $-9\%$ in adult earnings found by
\citet{page2009parental} in the PSID.\footnote{Comparisons of siblings are not used: siblings share and compete for the
family's resources and influence each other, and a job loss also affects later-born siblings through the family's
lasting income loss, so they are not a clean control. Code: \texttt{IGE/22.DisplacementID.py}.}

\begin{table}[htbp]\centering\small
\caption{Definitions of the father's job loss}\label{tab:displacement_definitions}
\resizebox{\textwidth}{!}{%
\begin{tabular}{lccc}
\toprule
 & \begin{tabular}[b]{@{}c@{}}(1) Log\\lifetime earnings\end{tabular} & \begin{tabular}[b]{@{}c@{}}(2) Log lifetime\\family income\end{tabular} & \begin{tabular}[b]{@{}c@{}}(3) Years\\of schooling\end{tabular} \\
\midrule
Closure or layoff (main) & $-0.092^{***}$ & $-0.079^{***}$ & $-0.431^{***}$ \\
 & (0.033) & (0.023) & (0.058) \\
Closures only & $-0.080^{*}$ & $-0.060^{**}$ & $-0.273^{***}$ \\
 & (0.042) & (0.030) & (0.077) \\
Quits (placebo) & $-0.065^{*}$ & $-0.020$ & $-0.296^{***}$ \\
 & (0.039) & (0.030) & (0.071) \\
Year without hours (earlier measure) & $-0.154^{***}$ & $-0.150^{***}$ & $-0.408^{***}$ \\
 & (0.047) & (0.034) & (0.075) \\
\bottomrule
\end{tabular}}
\begin{minipage}{0.95\textwidth}\footnotesize Exposure at child ages 0--17, all children; controls as in Table \ref{tab:displacement_children}. *** $p<0.01$, ** $p<0.05$, * $p<0.1$.\end{minipage}
\end{table}


\paragraph{What the model takes from this.} The PSID associations are an upper bound on the total effect of a father's
job loss on his children. They are robust only for White children's schooling and for job losses during childhood.
Administrative data find smaller effects: about $-9\%$ in adult earnings in Canada
\citep{oreopoulos2008intergenerational}, and much less in US tax data. There, a father's layoff during adolescence
cuts family income sharply, lowers college enrollment by less than half a percentage point, and leaves early-career
earnings almost unchanged \citep{hilger2016parental}. The model's child block maps family income at ages 0--5 and
6--17, and the share of childhood spent with both parents, into adult log earnings. The income slopes are set so that
a father's displacement at the mean displacement age lowers his child's adult earnings by 9\% (the high slope) or by
0.35\% (the low slope, from the US estimates of the working-paper version of \citet{hilger2016parental}, college quality measured by alumni earnings). Counterfactuals are reported under both. The presence slope comes from
the PSID. Holding family income at both childhood ages fixed, a childhood spent entirely with both parents is not
associated with higher adult earnings ($+0.011$, s.e.\ 0.045, $p=0.80$). It is associated with 0.30 more years of
schooling ($p<0.001$). Holding fixed only the family's income in the three years before and after the father's
displacement, the association with adult earnings is $0.19^{***}$. The model uses 0.19 as the high presence slope
and zero as the low one.\footnote{PSID children born 1968--1995 (SRC and SEO, unweighted) with at least three
childhood interviews. ``Both parents present'' means the head and the wife of the child's family unit are the child's
mother and father at birth. The share of childhood with both parents averages 0.38 for Black and 0.70 for White
children. Adult earnings are the mean over the years as head or wife at 25--55 (at least three years; zeros kept).
All regressions control for sex, birth year, parental education and, in pooled rows, race. Standard errors are
clustered by 1968 family; siblings are not compared. Without income controls, a childhood with both parents is
associated with adult earnings 29.5 log points higher ($0.295^{***}$, s.e.\ 0.041; $N=5{,}091$), with an earnings
rank 9.5 percentiles higher ($9.46^{***}$) and with 0.72 more years of schooling ($0.724^{***}$). Net of log family
income at 0--5 and 6--17, the earnings coefficient is $0.011$ (s.e.\ 0.045, $p=0.80$; $N=4{,}701$): Black $-0.031$
($p=0.65$), White $0.069$ ($p=0.24$). The schooling coefficient is $0.299^{***}$ (s.e.\ 0.081; $N=6{,}213$): Black
$0.213^{*}$ ($p=0.054$), White $0.476^{***}$. The income slopes in that regression are $0.148^{***}$ (ages 0--5) and
$0.320^{***}$ (6--17) for log earnings; they are associations, which is why the model takes its income slopes from
administrative estimates instead. Splitting the shares by age, the share with both parents at 6--17 is positive for earnings
($0.102^{**}$, $p=0.038$) and schooling ($0.342^{***}$). The share at 0--5 is negative for earnings ($-0.126^{**}$,
$p=0.032$) and not significant for schooling ($-0.080$, $p=0.41$). The two age shares are strongly collinear, and
their sum is close to the undivided coefficient. In the father's
displacement sample of Section~\ref{sec:children} (pooled, $N=3{,}520$), the regression of adult log earnings on the
displacement, the share with both parents, and log family income before and after the (pseudo) event gives
displacement $-0.040$ ($p=0.25$), presence $0.187^{***}$ (s.e.\ 0.054), income before $0.179^{***}$ and income after
$0.209^{***}$. Black and White estimates are $0.178^{**}$ ($p=0.037$) and $0.175^{**}$ ($p=0.013$). Without the
income after the event, the presence coefficient is $0.242^{***}$ and the displacement coefficient $-0.084^{**}$
($p=0.016$). Code: \texttt{IGE/79.ParentsPresentOutcomes.py}.} The parents' side of the model is estimated without
the child block, so its results for consumption and for the ranking of policies stand on their own.
